% Source: https://artofproblemsolving.com/community/c3h1930276_help_needed_in_problems_of_varied_topics
% Stable identity: aops-topic-1930276-q03
\textbf{Signs.} If $m<0$, then $0<3^m<1$ and $2^n=7+3^m$ lies strictly between $7$ and $8$; no power of $2$ does. So $m\ge0$, and then $2^n\ge8$, so $n\ge3$.

\textbf{$m=0$.} Then $2^n=8$ and $n=3$.

\textbf{$m\ge1$.} Since $n\ge3$, reduce modulo $8$: $3^m\equiv-7\equiv1\pmod8$. The powers of $3$ modulo $8$ alternate $3,1,3,1,\dots$, so $m$ is even, $m=2k$ with $k\ge1$. Modulo $3$: $2^n\equiv7\equiv1\pmod3$, and $2^n\equiv(-1)^n$, so $n$ is even, $n=2j$. Then
\[
7=2^{2j}-3^{2k}=\left(2^j-3^k\right)\left(2^j+3^k\right).
\]
Both factors are integers, the second is positive and larger, so $2^j-3^k=1$ and $2^j+3^k=7$. Hence $2^j=4$, $3^k=3$, i.e.\ $j=2$, $k=1$, giving $(m,n)=(2,4)$: $9+7=16$.

The solutions are $\boxed{(m,n)=(0,3)\text{ and }(2,4)}$.
